\chapter{Conclusion and Future Work}
\label{ch:conclusion}

\section{Conclusion}

A weather-adaptive ANPR pipeline has been built and tested up to real-time
plate reading. The YOLOv11n detector reached a mAP@50 of 99.2\,\% on more than
15\,000 plate images. EasyOCR read only 16.5--25.2\,\% of plates exactly in fog,
so it was replaced by LPRNet, whose character error fell from 0.70 to 0.249
after synthetic pre-training and training on real plates, and whose exact-plate
rate was raised further with public Indian plates, degradation training,
repeated-character plates and labelled real-video crops (38 to 41 of 51 plates
on six unseen videos). The real-time pipeline runs at up to 55 frames per second
on 1080p video and read 10 of 12 plates exactly on an unseen 4K traffic video.

For restoration, CLAHE was the fastest and most dependable of the first three
methods, and a Mamba-based network (UVMNet), trained in stages on traffic haze,
plate crops and real-fog pairs, gave much cleaner fog-free pictures. A
physics-guided procedure that uses the network only to measure the fog gives
the most natural pictures. A study of 60 real photographs taken at Ponmudi
showed, however, that none of these dehazers raised the number of plates read
beyond the level of a weather-trained reader on the untouched photographs, and
that combining several frames of the same vehicle is the effective lever. A
lightweight deraining network reached 29.19~dB on synthetic rain, but its effect
on real rain and on plate reading is not yet established.

The remaining errors are shortened repeated digits, plates from other states,
and small, blurred or faded plates. The vehicle database, alerts, dashboard,
dispatch and multi-node correlation are still to be built.

\section{Future Work}

\begin{itemize}
  \item Build the central system of Phase~2: vehicle database and status
        classification, real-time dashboard and alerts, dispatch by
        road-network distance, and multi-node correlation
        (Chapter~\ref{ch:phase2}).
  \item Port the node to a Jetson-class module and measure latency, memory and
        power; decide where the UVMNet-based restorer runs.
  \item Evaluate deraining and the weather router on real rain and night video,
        comparing plate reading with and without restoration.
  \item Collect and label more footage in fog, rain and at night, and use
        higher-resolution cameras for two-wheeler and auto plates, which are
        too small to read at present.
  \item Use multi-frame fusion in the live loop for vehicles whose vote share is
        low, and extend the format decoder to more state series.
\end{itemize}
